\section{Experiments}
\label{sec:experiments}

\subsection{Experimental Setup}
\label{sec:setup}

\looseness=-1
We apply all three analyses (\S\ref{sec:cross_arch_analyses}) to the $3 \times 4$ architecture-suite matrix (10 of 12 cells filled; Figure~\ref{fig:result_matrix}, Table~\ref{tab:cross_arch_matrix}).
For trained models, counterfactual sensitivity uses Cohen's $d$ on action MSE; reference VLMs use cosine distance on hidden representations (different scale; \S\ref{sec:cf_sensitivity}).
OpenVLA's OXE-pretrained checkpoint produces actions, permitting same-metric comparison (Bonferroni correction in Appendix~\ref{sec:bonferroni_correction}).

%------------------------------------------------------------------
\subsection{OpenVLA: Motivating Single-Architecture Analysis}
\label{sec:openvla_routing}

\looseness=-1
We first establish core phenomena on OpenVLA (Appendix~\ref{sec:openvla_deepdive_appendix}).
At L1, instruction-dominant heads form a gradient: random initialization 0/32, Llama-2 base 9/32, VLA-trained 24/32.
Per-layer $R^2$ probing (Figure~\ref{fig:action_r2_threecond}) shows the trained model peaks at L9 ($R^2 = 0.643$), Llama-2 base declines from L0, and random initialization collapses at deep layers.
All three share L0 ($R^2 = 0.645$); divergence arises from Transformer processing, with VLA training building on pretrained scaffold rather than constructing representations from scratch.

\begin{figure}[t]
\centering
\includegraphics[width=0.64\columnwidth]{figures/fig_action_r2_threecond.pdf}
\caption{Per-layer action-probe $R^2$ for three OpenVLA conditions (3 seeds, mean $\pm$ std); all share L0 ($R^2 = 0.645$). Trained model peaks at L9, Llama-2 base declines from L0, random initialization collapses at deep layers.}
\label{fig:action_r2_threecond}
\end{figure}

\begin{figure}[t]
\centering
\includegraphics[width=\columnwidth]{figures/fig_counterfactual_bar.pdf}
\caption{Counterfactual sensitivity, measured as Cohen's $d$ (OpenVLA, TraceVLA; dashed: $d{=}0.8$) and cosine distance (NORA, ref.\ VLMs; \S\ref{sec:cf_repr}). E=empty, W=wrong, S=shuffled, P=paraphrased.}
\label{fig:counterfactual_bar}
\end{figure}

%------------------------------------------------------------------
\subsection{Exit Collapse: Layer-Wise Action Probing}
\label{sec:r2_collapse}

\looseness=-1
We define $R^2$ exit collapse as $\Delta R^2 = R^2_{\text{peak}} - R^2_{\text{exit}}$ (Table~\ref{tab:cross_arch_matrix}).
VLM-to-VLA fine-tuning shows substantial exit collapse: TraceVLA $\Delta R^2 = 0.156$--$0.362$ across all four suites, NORA $\Delta R^2 = 0.308$--$0.455$ (Figure~\ref{fig:r2_collapse}), all exceeding the $0.10$ threshold.
Under VLA-to-VLA conditions, OpenVLA shows only small fluctuations ($\Delta R^2 = 0.033$--$0.086$); references similarly show $\Delta R^2 \leq 0.097$.
Within each architecture, collapse magnitude varies across suites, with TraceVLA ranging from $0.156$ (Spatial) to $0.362$ (Object), suggesting that task demands modulate the degree of exit-layer specialization, though all six VLM-to-VLA cells exceed the $0.10$ threshold.
A permutation test rejects random grouping ($p = 0.008$; Appendix~\ref{sec:delta_r2_threshold}), and Monte Carlo cross-validation confirms the separation: VLM-to-VLA trained CIs lie entirely above zero, while VLA-to-VLA CIs span zero (Table~\ref{tab:cross_arch_matrix}).

\begin{figure}[t]
\centering
\includegraphics[width=\columnwidth]{figures/fig_r2_collapse.pdf}
\caption{\textbf{Layer-wise action-probe $R^2$ (LIBERO-Goal, readout token).} Solid blue: trained; dashed orange: reference; dotted gray: NLP finetuning control ($R^2 = 0.155$, matching backbone family). VLM-to-VLA transitions exhibit pronounced peak-to-exit drops ($\Delta R^2 = 0.16$--$0.46$); VLA-to-VLA remains flat ($< 0.09$); NLP controls produce flat profiles ($\Delta R^2 = 0.000$). Demo-level splitting; see Table~\ref{tab:cross_arch_matrix}.}
\label{fig:r2_collapse}
\end{figure}

\paragraph{NLP finetuning control.}\looseness=-1
To test whether exit collapse is specific to action-oriented training, we probe three chat-finetuned models sharing backbones with our VLAs but lacking action heads or robot data: Llama-2-7b-chat, Phi-3-mini-4k-instruct, and Qwen2.5-3B-Instruct.\footnote{Qwen2.5-3B-Instruct shares the Qwen2.5 language backbone with NORA's Qwen2.5-VL-3B but lacks the vision encoder; see Appendix~\ref{sec:nlp_control_appendix}.}
All three produce flat $R^2$ profiles ($\Delta R^2 = 0.000$; dashed gray lines in Figure~\ref{fig:r2_collapse}), contrasting with VLM-to-VLA collapse ($\Delta R^2 = 0.156$--$0.455$).
\textbf{Exit collapse is consistently associated with VLM-to-VLA transitions and absent in both VLA-to-VLA and NLP-only fine-tuning.}

\begin{figure}[t]
\centering
\includegraphics[width=0.50\columnwidth]{figures/fig_multi_layer_attn.pdf}\hfill
\includegraphics[width=0.42\columnwidth]{figures/fig_dose_response.pdf}
\caption{\textbf{Left:} Instruction-dominant heads across all 32 layers (L1, L31 entry-exit for VLA-trained; L16 peak for Llama-2 base; zero for random). \textbf{Right:} Dose-response ablation (progressively masking instruction-dominant heads); divergence rises through 12 heads then plateaus. 95\% bootstrap CIs.}
\label{fig:synthesis_panels}
\end{figure}

\begin{table}[t]
\centering
\caption{\textbf{Peak-to-exit action-probe drop across the architecture-suite matrix.} $\Delta R^2$ is computed as the difference between the maximum layer-wise action-probe $R^2$ and the exit-layer $R^2$. Values above 0.10 are treated as substantial descriptive drops. OpenVLA is a VLA-to-VLA comparison, whereas TraceVLA and NORA are VLM-to-VLA comparisons.}
\label{tab:cross_arch_matrix}
\small
\setlength{\tabcolsep}{3pt}
\begin{tabular}{@{}lll cccc@{}}
\toprule
& & & \multicolumn{4}{c}{\textbf{$\Delta R^2$ (peak $-$ exit)}} \\
\cmidrule(l){4-7}
\textbf{Architecture} & \textbf{Suite} & \textbf{Transition} & \textbf{Trained} & \textbf{95\% CI} & \textbf{Reference} & \textbf{95\% CI} \\
\midrule
OpenVLA & Goal    & VLA$\to$VLA & 0.086 & [$-$0.060, 0.110] & 0.056 & [0.010, 0.080] \\
OpenVLA & Spatial & VLA$\to$VLA & 0.037 & [0.010, 0.100]    & 0.044 & [0.030, 0.110] \\
OpenVLA & Object  & VLA$\to$VLA & 0.037 & [0.000, 0.090]    & 0.051 & [0.000, 0.090] \\
OpenVLA & Long    & VLA$\to$VLA & 0.033 & [0.000, 0.100]    & 0.035 & [0.030, 0.100] \\
\midrule
TraceVLA & Goal    & VLM$\to$VLA & 0.333 & [0.210, 0.450]    & 0.065 & [$-$0.110, 0.050] \\
TraceVLA & Spatial & VLM$\to$VLA & 0.156 & [0.090, 0.220]    & 0.048 & ---$^{a}$ \\
TraceVLA & Object  & VLM$\to$VLA & 0.362 & [0.280, 0.500]    & 0.009 & [$-$0.070, 0.020] \\
TraceVLA & Long    & VLM$\to$VLA & 0.358 & [0.170, 0.530]    & 0.097 & [$-$0.020, 0.110] \\
\midrule
NORA    & Goal    & VLM$\to$VLA & 0.455$^{b}$ & ---$^{b}$  & 0.081 & [$-$0.010, 0.110] \\
NORA    & Spatial & VLM$\to$VLA & 0.308 & [0.300, 0.510]    & 0.089 & [$-$0.010, 0.090] \\
\bottomrule
\multicolumn{7}{l}{\footnotesize NORA Object and Long-Horizon checkpoints are unavailable.}\\
\multicolumn{7}{l}{\footnotesize $^{a}$CI artifact from flat $R^2$ profile. $^{b}$Bootstrap CI unreliable (exit-layer $R^2 \approx -0.006$).}\\
\end{tabular}
\end{table}

%------------------------------------------------------------------
\subsection{Instruction Sensitivity: Counterfactual Perturbation}
\label{sec:cf_sensitivity}

\paragraph{Instruction Sensitivity in Action-Producing VLAs}\looseness=-1
Across OpenVLA and TraceVLA, all perturbations yield $d = 0.45$--$1.52$ (Figure~\ref{fig:counterfactual_bar}; summary in Appendix Table~\ref{tab:cf_summary}, full rankings in Table~\ref{tab:counterfactual}).
``Empty'' ranks first in 6/8 cells; ``paraphrased'' ranks last in 7/8, indicating encoding beyond token-level matching.
TraceVLA produces systematically higher effect sizes ($d = 0.97$--$1.52$) than OpenVLA ($0.45$--$1.21$).
The ranking reveals a semantic gradient: paraphrased instructions retain full semantic content but alter surface form, yielding the smallest effects; empty strings remove all linguistic signal, producing the largest; wrong-task and shuffled-token conditions fall in between, indicating sensitivity to both task identity and token ordering.
TraceVLA's larger effect sizes are consistent with its spatial trace conditioning, which grounds instructions in visual trajectory tokens and may amplify instruction-action coupling.
\textbf{Instruction perturbation produces graded action divergence with effect sizes reflecting perturbation severity.}

\paragraph{Representation-Space Sensitivity in NORA and Reference VLMs}
\label{sec:cf_repr}

NORA and reference VLMs lack action outputs; their sensitivity is measured via cosine distance on hidden representations (different scale from Cohen's $d$; Appendix~\ref{sec:cf_details_appendix}).
The trained-reference contrast is consistent with action-space findings: NORA Goal saturates ($1.0$); reference VLMs show weak sensitivity ($\leq 0.19$).

%------------------------------------------------------------------
\subsection{Inherited Routing: L1 Attention Profiling}
\label{sec:l1_priors}

\looseness=-2
Counterfactual sensitivity confirms trained VLAs use instruction content, but not whether this reliance originates from training or backbone priors.
L1 attention profiling compares head routing before and after fine-tuning: TraceVLA's reference saturates at 32/32 heads; training reduces to 27--30/32.
NORA's reference allocates 13--14/16; training largely preserves this (13--14$\to$12--13/16).
OpenVLA shows 24--26/32 in the OXE-pretrained baseline with similar trained values (Appendix Figure~\ref{fig:cross_arch_l1}).
\textbf{L1 instruction routing is an inherited backbone property, not a consequence of VLA training.}

%------------------------------------------------------------------

\subsection{Synthesis: Training-Associated vs.\ Inherited Routing}
\label{sec:routing_synthesis}

\looseness=-2
The three analyses separate two patterns: exit collapse and counterfactual sensitivity are training-associated (pronounced in VLM-to-VLA, small in VLA-to-VLA), while L1 attention is inherited from reference backbones.
Trained OpenVLA and TraceVLA converge to nearly identical mid-layer peaks ($R^2 = 0.637$ at L9 vs.\ $0.634$ at L11; Appendix~\ref{sec:cross_arch_r2}), despite differing backbone families, suggesting a shared representational bottleneck in autoregressive VLA fine-tuning.
The 32-layer profile (Figure~\ref{fig:synthesis_panels}, left) shows VLA-trained entry-exit concentration (L1: 24, L31: 18) versus Llama-2 base mid-depth peak (L16: 25); dose-response ablation (right) yields monotonic divergence ($d = 0.48$--$0.64$), causally linking L1 gating to action computation (Appendix Table~\ref{tab:specificity_control}), confirmed across all four LIBERO suites and absent in three NLP controls (Appendix Table~\ref{tab:convergent_evidence}).
